\chapter{Conclusion}

\section{Introduction}

This chapter summarises the key findings of the study, discusses theoretical and practical implications, and outlines limitations and directions for future research.

\section{Summary of Findings}

This dissertation examined how artificial intelligence (AI) and generative AI influence consumer perception in marketing. Through a thematic analysis of existing literature, five key themes were identified: trust, authenticity, personalisation, privacy concerns, and purchase intention.

The findings suggest that AI has a dual impact on consumer perception. On one hand, AI improves efficiency, relevance, and personalisation in marketing. On the other hand, it raises concerns related to authenticity, transparency, and privacy. These conflicting effects indicate that AI-driven marketing involves a trade-off between benefits and risks.

Trust was found to be a central factor linking all other themes. Consumers are more likely to accept AI-driven marketing when they perceive it as trustworthy and useful. However, when authenticity is questioned or privacy concerns arise, trust may decrease, leading to lower purchase intention.

\section{Theoretical Implications}

This study contributes to the literature by integrating research on AI in marketing with consumer perception theories. It highlights the importance of examining multiple interconnected factors rather than focusing on a single variable.

In particular, the study emphasises the role of authenticity and trust in shaping consumer responses to AI-generated content. It also extends existing research by considering generative AI as a distinct category that directly influences marketing communication.

\section{Practical Implications}

From a practical perspective, the findings suggest that marketers should carefully balance the use of AI technologies with consumer expectations.

First, transparency is important. Clearly communicating the use of AI may help reduce uncertainty and increase trust.

Second, firms should ensure that AI-generated content maintains a level of authenticity, especially in emotional or human-centred communication.

Third, data usage should be managed responsibly to reduce privacy concerns and build long-term consumer trust.

\section{Limitations}

This study has several limitations. It relies on secondary data and does not include primary empirical research. The selection of literature may also introduce bias, as only a limited number of studies were analysed.

In addition, the findings are based on interpretation and may not fully capture the diversity of consumer experiences across different cultural and market contexts.

\section{Future Research}

Future research could address these limitations by conducting empirical studies, such as surveys or experiments, to test the relationships identified in this study.

Further research is also needed to explore how different types of AI applications influence consumer perception in various contexts. In particular, the impact of generative AI on authenticity and trust remains an important area for investigation.

\section{Final Remarks}

Artificial intelligence and generative AI are transforming the marketing landscape. While they offer significant opportunities for improving efficiency and personalisation, they also introduce new challenges for consumer perception.

Understanding how consumers respond to AI-driven marketing is therefore essential for both researchers and practitioners. This study provides a foundation for future research in this evolving field.